\section{Grammar Keystone 240 Model}

\subsection{Purpose}

The breakthrough portfolio ranked \texttt{240} as the highest-leverage target. This section executes that target. It tests whether \texttt{240} should be treated as a reusable grammar operator, an ordinary stem, a formula boundary sign, a terminal marker, or a mixed sign. The model is still pre-phonetic: it licenses controlled environments, not sound values.

\subsection{Summary}

\begin{table}[htbp]
\centering
\footnotesize
\begin{tabular}{p{0.30\textwidth}p{0.24\textwidth}p{0.32\textwidth}}
\toprule
\textbf{Metric} & \textbf{Value} & \textbf{Note} \\
\midrule
240 token occurrences & 354 & texts=354 \\
Position profile & Medial:234; Initial:96; Final:16; Standalone:8 & Reading-order positions in cleaned corpus. \\
Distinct signs after 240 & 69 & 520:46; 740:25; 233:24; 636:19; 482:15 \\
Distinct signs before 240 & 28 & 002:83; 235:45; 798:20; 060:16; 220:15 \\
Observed local frames & 182 & 002|520:29; 235|740:15; $<$START$>$|636:10; $<$START$>$|705:9; $<$START$>$|233:9 \\
240+X slot groups & 54 & 18 frame-locked groups at score $>$= 0.70 \\
Controlled 240+STEM lanes & 3 & Can feed later phonetic tests after image checks. \\
Validated 240 morphology constraints & 8 & Promoted to grammar model only. \\
Leading 240 model & Qualifier/title/classifier prefix operator & score=0.986; status=LEADING MODEL \\
\bottomrule
\end{tabular}
\caption{Summary of the \texttt{240} grammar-keystone test.}
\label{tab:grammar-keystone-240-summary}
\end{table}

\subsection{Rival Models}

\begin{table}[htbp]
\centering
\footnotesize
\begin{tabular}{p{0.52\textwidth}rp{0.22\textwidth}}
\toprule
\textbf{Model} & \textbf{Score} & \textbf{Status} \\
\midrule
Qualifier/title/classifier prefix operator & 0.986 & LEADING MODEL \\
Polyfunctional mixed sign & 0.615 & CONTROL RISK \\
Formula opener/boundary marker & 0.463 & SECONDARY FUNCTION \\
Terminal/suffix/numeral value & 0.060 & LOW SUPPORT \\
Ordinary lexical stem & 0.012 & WEAKENED \\
\bottomrule
\end{tabular}
\caption{Rival grammatical models for \texttt{240}.}
\label{tab:grammar-keystone-240-models}
\end{table}

\subsection{240+X Families}

\begin{table}[htbp]
\centering
\scriptsize
\begin{tabular}{lrlrp{0.40\textwidth}}
\toprule
\textbf{Post240} & \textbf{Count} & \textbf{Top frame} & \textbf{Lock} & \textbf{Lead} \\
\midrule
520 & 33 & 002|$<$END$>$ & 0.800 & Prime 240+X environment; control grammar before treating X as a stem. \\
$<$NONE$>$ & 31 & $<$START$>$|233 & 0.476 & 240 occurs as a complete filler; treat as possible standalone operator/control. \\
740 & 21 & 235|$<$END$>$ & 0.843 & Frame-locked 240 compound; useful grammar control. \\
482 & 16 & 002|740 & 0.933 & Prime 240+STEM environment; use as controlled stem lane. \\
415 & 12 & 002|740 & 0.817 & Mixed frame evidence; keep as lower-confidence 240 context. \\
636 & 11 & 031|$<$END$>$ & 0.579 & Mixed frame evidence; keep as lower-confidence 240 context. \\
220 & 7 & 002|$<$END$>$ & 0.739 & Mixed frame evidence; keep as lower-confidence 240 context. \\
642 & 7 & 060|$<$END$>$ & 0.436 & Mixed frame evidence; keep as lower-confidence 240 context. \\
904 & 7 & 002|740 & 0.751 & Prime 240+STEM environment; use as controlled stem lane. \\
603 & 6 & 060|$<$END$>$ & 0.533 & Mixed frame evidence; keep as lower-confidence 240 context. \\
\bottomrule
\end{tabular}
\caption{Highest-yield \texttt{240+X} families.}
\label{tab:grammar-keystone-240-families}
\end{table}

\subsection{Frame Identity Tests}

\begin{table}[htbp]
\centering
\scriptsize
\begin{tabular}{lrrrp{0.38\textwidth}}
\toprule
\textbf{Post240} & \textbf{With} & \textbf{Bare} & \textbf{Overlap} & \textbf{Result} \\
\midrule
176 & 4 & 14 & 0.571 & 240 stem lane \\
482 & 16 & 3 & 0.333 & 240 stem lane \\
520 & 33 & 42 & 0.132 & 240 grammar control \\
904 & 7 & 7 & 0.286 & 240 stem lane \\
220 & 7 & 23 & 0.261 & mixed/low evidence \\
740 & 21 & 34 & 0.261 & frame-locked; needs bare control \\
806 & 2 & 10 & 0.300 & mixed/low evidence \\
235 & 3 & 35 & 0.143 & mixed/low evidence \\
773 & 5 & 1 & 0.000 & mixed/low evidence \\
233 & 4 & 13 & 0.481 & mixed/low evidence \\
\bottomrule
\end{tabular}
\caption{Tests of whether \texttt{240+X} preserves enough frame identity to become a controlled environment.}
\label{tab:grammar-keystone-240-preservation}
\end{table}

\subsection{Decision}

The leading model treats \texttt{240} as a qualifier/title/classifier-like operator. This does not decipher \texttt{240}; it constrains it. The working notation should therefore keep \texttt{240} as an abstract operator, \texttt{OP240}, while the following position is tested as the stem-like slot.

The practical consequence is large: \texttt{240-482}, \texttt{240-904}, \texttt{240-176}, and related forms can now be tested as one grammar family. The first phonetic lane remains conditional on image review, but it no longer floats without structure.

\subsection{Acceleration Queue}

\begin{table}[htbp]
\centering
\scriptsize
\begin{tabular}{rp{0.34\textwidth}p{0.38\textwidth}}
\toprule
\textbf{Priority} & \textbf{Action} & \textbf{Output} \\
\midrule
1 & Freeze 240 as a grammar operator in the solver & Update constrained readings so 240 is an operator label, not a sound-bearing value. \\
2 & Promote the first phonetic-grade 240+STEM lane & Image-check and context-check 240-482 against 240-904 as one controlled contrast. \\
3 & Use 176 as the bridge/control stem & Compare bare 176, 240-176, 048, and 061 inside the prime entity frame. \\
4 & Use frame-locked 240 compounds as controls & Generate a 240+STEM contrast matrix for 482, 904, 176, 773, and high-count followers. \\
5 & Split 240 from 235 before lexical claims & Build a rival-opener model: 240+X versus 235+X and X+terminal contexts. \\
6 & Block sound values for 240 & Keep 240 as OP240 in abstract reading skeletons. \\
\bottomrule
\end{tabular}
\caption{Immediate actions unlocked by the \texttt{240} grammar model.}
\label{tab:grammar-keystone-240-actions}
\end{table}

\subsection{Outputs}

\begin{itemize}
  \item \texttt{outputs/grammar\_keystone\_240\_summary.csv}
  \item \texttt{outputs/grammar\_keystone\_240\_occurrences.csv}
  \item \texttt{outputs/grammar\_keystone\_240\_compound\_families.csv}
  \item \texttt{outputs/grammar\_keystone\_240\_preservation\_tests.csv}
  \item \texttt{outputs/grammar\_keystone\_240\_operator\_models.csv}
  \item \texttt{outputs/grammar\_keystone\_240\_action\_queue.csv}
\end{itemize}
